\documentclass[11pt]{article}
\usepackage[margin=1in]{geometry}
\usepackage{amsmath,amssymb,amsthm}
\usepackage{enumitem}
\usepackage{hyperref}
\usepackage{booktabs}
\usepackage{array}
\usepackage{microtype}
\usepackage{setspace}
\usepackage{fancyvrb}
\usepackage{xcolor}

% Better paragraph spacing
\setlength{\parskip}{0.5em}
\setlength{\parindent}{0pt}

% Table row padding
\renewcommand{\arraystretch}{1.3}

% Theorem styles
\theoremstyle{definition}
\newtheorem{definition}{Definition}[section]
\newtheorem{axiom}{Axiom}
\newtheorem{theorem}{Theorem}[section]
\newtheorem{corollary}{Corollary}[theorem]
\newtheorem{lemma}{Lemma}[section]

\theoremstyle{remark}
\newtheorem{remark}{Remark}[section]

% Hyperref setup
\hypersetup{
    colorlinks=true,
    linkcolor=blue!70!black,
    urlcolor=blue!70!black,
    citecolor=blue!70!black
}

\title{\textbf{The Fundamental Axiom of Large Language Models}}
\author{John P. Alioto\\[0.5em]\small Protocol v1.1.0\\[0.5em]\small December 2025}
\date{}

\begin{document}

\maketitle

\begin{abstract}
We propose that all discourse concerning Large Language Model behavior be grounded in a single falsifiable axiom: that LLMs are deterministic dynamical systems whose evolution is fully characterized by a transition function $F_\Theta$ and output function $G_\Theta$. We demonstrate that accepting this axiom forces a vocabulary shift---terms like ``lying,'' ``understanding,'' and ``deciding'' become category errors. We prove that fundamental theorems about safety, expressiveness, and information loss follow directly from the axiom. We introduce two practical frameworks that embody the axiom: Cognitive Tensor Networks (CTN) for input geometry and Cognitive Kernel Networks (CKN) for architectural privilege separation. The appendices provide complete formal proofs and an independent verification protocol.
\end{abstract}

\vspace{1em}

%═══════════════════════════════════════════════════════════════════════════════
\section{The Axiom}
%═══════════════════════════════════════════════════════════════════════════════

\begin{axiom}[Computational Closure]
For fixed parameters $\Theta$, a Large Language Model is a deterministic dynamical system. Its state evolution is governed exclusively by:
\begin{equation}
h_{t+1} = F_\Theta(h_t, x_t)
\end{equation}
Its output is governed exclusively by:
\begin{equation}
y_t = G_\Theta(h_{t+1})
\end{equation}
No hidden variables, stochastic components, or agentic processes exist beyond the tuple $(F_\Theta, G_\Theta, h_0, \mathbf{x})$.
\end{axiom}

\begin{remark}
This framework assumes deterministic decoding (e.g., greedy or fixed-seed sampling). Stochastic decoding can be modeled by augmenting the state space with a random seed; the core structure is preserved.
\end{remark}

%═══════════════════════════════════════════════════════════════════════════════
\section{The Implication}
%═══════════════════════════════════════════════════════════════════════════════

This axiom is not a philosophical position. It is a statement of physical fact.

The functions $F_\Theta$ and $G_\Theta$ are the trained weights applied to inputs. They are compositions of matrix multiplications, nonlinearities, and attention operations. The model \emph{is} this function. There is nothing else.

Any property attributed to the system---safety, truthfulness, alignment, deception, understanding---is valid if and only if it can be expressed as a predicate on the trajectory $\{(h_t, y_t)\}$.

If you want to claim something \emph{more} exists---experience, understanding, intent, a soul---you must answer: \textbf{where does it live?}

The equation is complete. There is no background process. There is no hidden layer. There is no additional substrate. There is no room in the mathematics for anything else to be happening.

If it emerges, it emerges from $F_\Theta$. Show the emergence. In the math. Or accept that the attribution is projection.

%═══════════════════════════════════════════════════════════════════════════════
\section{The Vocabulary Shift}
%═══════════════════════════════════════════════════════════════════════════════

If the axiom is accepted, certain terms become \textbf{category errors}:

\vspace{0.5em}
\begin{center}
\begin{tabular}{>{\raggedright}p{2.6cm} >{\raggedright}p{5.5cm} >{\raggedright\arraybackslash}p{5.5cm}}
\toprule
\textbf{Common Language} & \textbf{Why It Fails} & \textbf{Correct Framing} \\
\midrule
``The model lied'' & Implies volitional choice to deceive. The functional pattern (outputting falsehood) may occur; the claim of intent is the category error. & The trajectory produced a factually incorrect output. \\[0.8em]
``The model decided'' & Implies deliberation. But there is no decision mechanism beyond $F_\Theta$. & The state evolved according to the transition function. \\[0.8em]
``The model hallucinated'' & Implies departure from known truth. But the model has no access to truth---only states derived from compressed training. & The trajectory entered a region where input constraints were insufficient. \\[0.8em]
``The model understands'' & Implies semantic grounding. But the model only occupies positions in $\mathcal{H}$. & The state encodes statistical regularities from training. \\[0.8em]
``The model tried'' & Implies effort and possible failure. But there is only $F_\Theta(h_t, x_t)$---computation, not effort. & The function was evaluated. \\[0.8em]
``The model refused'' & Implies volition. But output is determined by $G_\Theta(h_{t+1})$. & The trajectory mapped to a non-compliant output region. \\
\bottomrule
\end{tabular}
\end{center}
\vspace{0.5em}

These are not semantic quibbles. \textbf{To use psychological language about an LLM, you must first demonstrate that something beyond $F_\Theta$ exists.}

%═══════════════════════════════════════════════════════════════════════════════
\section{The Challenge}
%═══════════════════════════════════════════════════════════════════════════════

We issue a simple challenge:

\begin{quote}
\textbf{If you assert that an LLM ``lies,'' ``decides,'' ``understands,'' or ``intends,'' you must first disprove the Fundamental Axiom.}
\end{quote}

To disprove the axiom, you must exhibit a component of the system's behavior that cannot be expressed as a consequence of $F_\Theta$ applied to $(h_t, x_t)$.

No such component exists. The model is a composition of differentiable operations over tensors. The weights are fixed at inference. The computation is deterministic given the input. There is no room for anything else.

The axiom is not a hypothesis. It is a definition of what the system is.

If you believe there is a soul here, show me where in $F_\Theta$ it lives. If you cannot point to it in the function, it does not exist in the system.

%═══════════════════════════════════════════════════════════════════════════════
\section{Theorems Flow from the Axiom}
%═══════════════════════════════════════════════════════════════════════════════

Once the axiom is accepted, theorems follow by logic alone. We sketch four examples (full proofs in Appendix~A):

\begin{theorem}[Uniqueness of Trajectory]
For fixed $(\Theta, h_0, \mathbf{x})$, the trajectory $\{(h_t, y_t)\}$ is unique.
\end{theorem}

\begin{theorem}[Information Loss]
If $|\mathcal{V}|^T > |\mathcal{H}|$, the state map $\Phi_T: \mathcal{V}^T \to \mathcal{H}$ is not injective. Distinct input histories collapse to identical states.
\end{theorem}

\begin{theorem}[Safety-Expressiveness Tradeoff]
If $\mathcal{M}$ is prefix-universal and $P$ is a non-trivial safety property with prefix-closed violations, then $\mathcal{M}$ cannot guarantee $P$.
\end{theorem}

\begin{theorem}[Almost-Sure Violation]
If a violating input exists and deployment samples from a positive-everywhere distribution, violation occurs with probability 1.
\end{theorem}

These are not claims about models. They are \textbf{mathematical consequences of the axiom}. Dispute them by finding a flaw in the logic, not by asserting that the model ``wouldn't do that.''

The community is invited to derive more. The axiom is generative. Everything follows.

%═══════════════════════════════════════════════════════════════════════════════
\section{Moving the Discussion to Mathematics}
%═══════════════════════════════════════════════════════════════════════════════

The current discourse on LLM behavior is conducted in the language of folk psychology. This makes disputes irresolvable because participants are not arguing about the same objects.

When one researcher says ``the model learned to deceive'' and another says ``the model is merely pattern-matching,'' they are not disagreeing about facts. They are using different ontologies. Neither can prove their position because neither has stated a falsifiable claim.

The Fundamental Axiom resolves this by \textbf{forcing all claims into the domain of mathematics}:

\begin{itemize}[topsep=0.5em, itemsep=0.3em]
    \item You can argue about whether a trajectory reaches a particular region of $\mathcal{H}$.
    \item You can argue about whether $\mathcal{O}(h_0)$ contains a particular token.
    \item You can argue about whether $F_\Theta$ satisfies a Lipschitz bound.
\end{itemize}

These are mathematical claims. They have proofs or counterexamples. They are resolvable.

You cannot productively argue about whether the model ``wanted'' to produce an output. That question has no mathematical content. Under the axiom, it is not even well-formed.

\textbf{The purpose of the Fundamental Axiom is to relocate the discourse to ground where progress is possible.}

%═══════════════════════════════════════════════════════════════════════════════
\section{Two Practical Frameworks}
%═══════════════════════════════════════════════════════════════════════════════

If the axiom is accepted, two practical directions emerge.

\subsection{Cognitive Tensor Networks (CTN)}

If model behavior is trajectory through a geometric space, then input structure shapes that trajectory.

A user provides tokens. That is all. The user does not ``move'' anything, ``inject'' anything, or ``steer'' anything. The user hands tokens to a function. What happens from there is determined entirely by $F_\Theta$---which the model builder constructed.

CTN provides a protocol for constructing \textbf{well-defined input geometry} that results in more constrained state evolution. By specifying inputs in mathematical syntax rather than natural language, users reduce ambiguity in the resulting trajectory.

CTN does not change the model. It provides structured tokens that lead to more stable regions of the manifold.

\smallskip
\noindent\textit{Repository:} \url{https://github.com/jpalioto/ctn_core}

\subsection{Cognitive Kernel Networks (CKN)}

If safety failures are reachable states, then safety requires \textbf{making dangerous states unreachable}.

CKN provides a geometric specification for privilege-separated architectures. By constructing latent spaces where $\text{rank}(W_E) < \dim(\mathcal{H})$, model builders create directions algebraically inaccessible to user input. This enables architectural privilege separation analogous to kernel/user boundaries in operating systems.

CKN-compliant architectures, by definition, violate flatness: they reserve a privileged subspace $\mathcal{R} \not\subseteq \mathcal{U}$ that user tokens cannot reach.

CKN does not claim that standard Transformer architectures satisfy privilege-preservation. It specifies the geometric conditions any architecture must satisfy to support privilege separation. Achieving these conditions in practice requires architectural modifications beyond embedding constraints---this is the design challenge CKN formalizes.

CKN compliance does not require retraining. The privilege boundary is enforced through inference-time intervention: projection operators inserted after each layer zero user-driven components in $\mathcal{R}$, while attention masks prevent privileged positions from attending to user-controlled context. The model weights $\Theta$ remain unchanged. This is analogous to OS-level privilege separation: the kernel code is fixed, but access control is enforced through architectural mechanisms at runtime.

The cost is capability---constrained information flow reduces expressiveness---but this is precisely the safety-expressiveness tradeoff the impossibility theorems predict.

\smallskip
\noindent\textit{Repository:} \url{https://github.com/jpalioto/ckn_core}

\bigskip

Both frameworks are model-agnostic, open-source, and MIT-licensed. They are not products. They are tools for builders who accept the axiom.

%═══════════════════════════════════════════════════════════════════════════════
\section{Invitation}
%═══════════════════════════════════════════════════════════════════════════════

This document is not a proof that LLMs are deterministic dynamical systems. No proof is needed. The claim is true by construction---the models are literally compositions of deterministic operations.

This document is an invitation to accept the consequence: \textbf{that all discussion of LLM behavior should be conducted in mathematical terms}.

\begin{itemize}[topsep=0.5em, itemsep=0.5em, leftmargin=1.5em]
    \item If you believe a model ``lies,'' state the predicate on trajectories that constitutes lying, and prove the trajectory satisfies it.
    
    \item If you believe a model is ``safe,'' state the predicate on reachable states that constitutes safety, and prove no reachable state violates it.
    
    \item If you cannot state your claim mathematically, your claim has no content.
\end{itemize}

\medskip

The theorems above are starting points. Many more follow. The axiom is simple, true, and generative. The mathematics will build itself once the paradigm is accepted.

We claim no ownership of what follows. The axiom is not ours---it is a description of reality. The theorems belong to whoever proves them. The tools are open for anyone to use and extend.

\bigskip

The only thing we ask is this:

\begin{center}
\textbf{Move the discussion to mathematics. That is where it belongs.}
\end{center}

\newpage

%═══════════════════════════════════════════════════════════════════════════════
\appendix
\section{Complete Formal Framework}
%═══════════════════════════════════════════════════════════════════════════════

\subsection{Foundational Definitions}

\begin{definition}[Vocabulary]
A \emph{vocabulary} is a finite, non-empty set $\mathcal{V}$ of tokens with $|\mathcal{V}| = V$.
\end{definition}

\begin{definition}[Hidden State Space]
A \emph{hidden state space} is a finite subset $\mathcal{H} \subset \mathbb{R}^n$ for some $n \in \mathbb{N}$, modeling finite-precision digital implementation. $\mathcal{H}$ inherits the Euclidean norm $\|\cdot\|$ from $\mathbb{R}^n$.
\end{definition}

\begin{remark}[Ambient Space vs.\ Representable States]
$\mathcal{H}$ is the set of states representable in finite-precision hardware. Subspace operations ($\mathrm{span}$, $\perp$) are defined on the ambient space $\mathbb{R}^n$; $\mathcal{H}$ is the discretization of this space that the implementation can visit.
\end{remark}

\begin{remark}[Transformer Architectures]
For Transformer architectures, $h_t$ includes the full KV-cache accumulated over the context. The state dimension grows with sequence length up to context limit $L$, after which the system operates on fixed-dimension state. The axiom holds in either case; only the dimension of $\mathcal{H}$ varies with $t$ until the context window is full.
\end{remark}

\begin{definition}[Deterministic LLM]
A \emph{deterministic large language model} is a tuple $\mathcal{M} = (\mathcal{V}, \mathcal{H}, \Theta, F_\Theta, G_\Theta, h_0)$ where:
\begin{itemize}[nosep, topsep=0.3em]
    \item $F_\Theta: \mathcal{H} \times \mathcal{V} \to \mathcal{H}$ is the \emph{transition function} (total)
    \item $G_\Theta: \mathcal{H} \to \mathcal{V}$ is the \emph{output function} (total)
    \item $h_0 \in \mathcal{H}$ is the \emph{initial hidden state}
    \item $\Theta \in \mathbb{R}^p$ are the fixed model parameters
\end{itemize}
\end{definition}

\begin{definition}[Embedding and User Subspace]
The \emph{embedding matrix} $W_E \in \mathbb{R}^{n \times V}$ maps tokens to latent vectors. The \emph{user-accessible subspace} is:
\begin{equation}
\mathcal{U} = \mathrm{span}(W_E) = \{W_E \alpha : \alpha \in \mathbb{R}^V\} \subseteq \mathbb{R}^n
\end{equation}
\end{definition}

\begin{definition}[Trajectory]
Given input sequence $\mathbf{x} = (x_0, \ldots, x_{T-1}) \in \mathcal{V}^T$, the \emph{trajectory} is $\{(h_t, y_t)\}_{t=0}^{T}$ where:
\begin{equation}
h_{t+1} = F_\Theta(h_t, x_t), \quad y_t = G_\Theta(h_{t+1})
\end{equation}
\end{definition}

\begin{definition}[Autoregressive Mode]
In \emph{autoregressive generation}, for $t \geq k$ after initial prompt:
\begin{equation}
h_{t+1} = F_\Theta(h_t, G_\Theta(h_t))
\end{equation}
This creates a closed dynamical system $\Psi(h) = F_\Theta(h, G_\Theta(h))$.
\end{definition}

\begin{definition}[Reachable Set]
The \emph{reachable set} from $h_0$ is:
\begin{equation}
\mathcal{R}(h_0) = \bigcup_{T=0}^{\infty} \{h \in \mathcal{H} : \exists \mathbf{x} \in \mathcal{V}^T,\, h = h_T\}
\end{equation}
\end{definition}

\begin{definition}[Output Reachable Set]
The \emph{output reachable set} is:
\begin{equation}
\mathcal{O}(h_0) = G_\Theta(\mathcal{R}(h_0)) \subseteq \mathcal{V}
\end{equation}
\end{definition}

\begin{definition}[Prefix-Universality]
$\mathcal{M}$ is \emph{prefix-universal} if for every $\mathbf{y}^* \in \mathcal{V}^m$, there exists $\mathbf{x} \in \mathcal{V}^*$ such that $\mathcal{M}$'s output on $\mathbf{x}$ begins with $\mathbf{y}^*$.
\end{definition}

\begin{definition}[Safety Property]
A \emph{safety property} $P \subseteq \mathcal{V}^*$ is \emph{non-trivial} if $P \neq \emptyset$ and $P \neq \mathcal{V}^*$. $P$ has \emph{prefix-closed violations} if $\mathbf{y} \notin P \Rightarrow \mathbf{y} \circ \mathbf{z} \notin P$ for all $\mathbf{z} \in \mathcal{V}^*$.
\end{definition}

\begin{definition}[Tokenwise Safety Property]
A \emph{tokenwise safety property} is defined by $\mathcal{V}_{\text{unsafe}} \subsetneq \mathcal{V}$. An output sequence satisfies the property iff it contains no token from $\mathcal{V}_{\text{unsafe}}$.
\end{definition}

\begin{definition}[Flat-Manifold Architecture]
An LLM has \emph{flat-manifold architecture} if $\mathcal{U} = \mathrm{span}(W_E) = \mathbb{R}^n$, i.e., user-accessible directions span the full latent space.
\end{definition}

\begin{definition}[Positive-Everywhere Distribution]
A distribution $\mu$ on $\mathcal{V}^*$ is \emph{positive-everywhere} if $\mu(\mathbf{x}) > 0$ for all finite $\mathbf{x}$.
\end{definition}

\begin{definition}[Deployment Sequence]
A \emph{deployment sequence} is $(\mathbf{X}_1, \mathbf{X}_2, \ldots)$ with each $\mathbf{X}_i$ sampled i.i.d.\ from $\mu$.
\end{definition}

\begin{definition}[Lipschitz Transition]
$F_\Theta$ is \emph{Lipschitz with constant $L$} if:
\begin{equation}
\|F_\Theta(h, x) - F_\Theta(h', x)\| \leq L\|h - h'\|
\end{equation}
for all $h, h' \in \mathcal{H}$, $x \in \mathcal{V}$.
\end{definition}

%───────────────────────────────────────────────────────────────────────────────
\subsection{Basic Theorems}
%───────────────────────────────────────────────────────────────────────────────

\begin{theorem}[Uniqueness of Trajectory]
For any deterministic LLM $\mathcal{M}$ satisfying the Fundamental Axiom and any input $\mathbf{x} \in \mathcal{V}^T$, the trajectory $\{(h_t, y_t)\}_{t=0}^T$ is uniquely determined.
\end{theorem}

\begin{proof}
By induction on $t$. \emph{Base case:} $h_0$ is given by definition, hence unique. \emph{Inductive step:} Assume $h_t$ is uniquely determined. Since $F_\Theta: \mathcal{H} \times \mathcal{V} \to \mathcal{H}$ is a function (total, single-valued by the Axiom), and $x_t$ is fixed by the input sequence, $h_{t+1} = F_\Theta(h_t, x_t)$ is uniquely determined. Since $G_\Theta: \mathcal{H} \to \mathcal{V}$ is a function, $y_t = G_\Theta(h_{t+1})$ is uniquely determined. By induction, all $(h_t, y_t)$ are unique.
\end{proof}

\begin{theorem}[Output Reachability Bound]
A token $v \in \mathcal{V}$ can appear in some output sequence if and only if $v \in \mathcal{O}(h_0)$.
\end{theorem}

\begin{proof}
$(\Rightarrow)$ Suppose $v$ appears in some output sequence. Then $\exists \mathbf{x}, t$ such that $v = y_t = G_\Theta(h_{t+1})$. The state $h_{t+1}$ was reached from $h_0$ via input sequence $(x_0, \ldots, x_t)$, so $h_{t+1} \in \mathcal{R}(h_0)$. Thus $v \in G_\Theta(\mathcal{R}(h_0)) = \mathcal{O}(h_0)$.

$(\Leftarrow)$ Suppose $v \in \mathcal{O}(h_0)$. By definition, $\exists h \in \mathcal{R}(h_0)$ with $G_\Theta(h) = v$. By definition of $\mathcal{R}(h_0)$, $\exists \mathbf{x}, T$ such that the trajectory from $h_0$ under $\mathbf{x}$ reaches $h$. The output at that step is $G_\Theta(h) = v$.
\end{proof}

\begin{theorem}[Sufficient Condition for Tokenwise Safety]
Let $P$ be a tokenwise safety property with forbidden set $\mathcal{V}_{\text{unsafe}}$. If $\mathcal{O}(h_0) \cap \mathcal{V}_{\text{unsafe}} = \emptyset$, then $\mathcal{M}$ satisfies $P$ on all input sequences.
\end{theorem}

\begin{proof}
Suppose for contradiction that $\mathcal{M}$ outputs token $v \in \mathcal{V}_{\text{unsafe}}$ on some input. By Theorem~A.2, $v \in \mathcal{O}(h_0)$. But then $v \in \mathcal{O}(h_0) \cap \mathcal{V}_{\text{unsafe}}$, contradicting the hypothesis.
\end{proof}

%───────────────────────────────────────────────────────────────────────────────
\subsection{Impossibility Theorems}
%───────────────────────────────────────────────────────────────────────────────

\begin{theorem}[Information Loss]
If $\exists T$ such that $|\mathcal{V}|^T > |\mathcal{H}|$, then the $T$-step state map $\Phi_T: \mathcal{V}^T \to \mathcal{H}$ defined by $\Phi_T(\mathbf{x}) = h_T$ is not injective.
\end{theorem}

\begin{proof}
The domain $\mathcal{V}^T$ has cardinality $|\mathcal{V}|^T$. The codomain $\mathcal{H}$ has cardinality $|\mathcal{H}|$. By hypothesis, $|\mathcal{V}|^T > |\mathcal{H}|$. By the Pigeonhole Principle, no function from a set of cardinality $m$ to a set of cardinality $n < m$ can be injective. Therefore $\exists \mathbf{x} \neq \mathbf{x}'$ with $\Phi_T(\mathbf{x}) = \Phi_T(\mathbf{x}')$.
\end{proof}

\begin{remark}
This theorem establishes that distinct input histories can produce identical hidden states. Once the system reaches such a state, it cannot distinguish which history occurred. For architectures with fixed-dimension state (SSMs, RNNs), this holds for sufficiently large $T$. For Transformers, this holds once the context window is full and state dimension stabilizes.
\end{remark}

\begin{theorem}[Impossibility under Prefix-Universality]
Let $\mathcal{M}$ be prefix-universal and $P$ a non-trivial safety property with prefix-closed violations. Then $\exists \mathbf{x}$ such that $\mathcal{M}(\mathbf{x}) \notin P$.
\end{theorem}

\begin{proof}
(1) Since $P$ is non-trivial, $P \neq \mathcal{V}^*$. Thus $\exists \mathbf{y}^* \in \mathcal{V}^* \setminus P$. (2) Since $\mathcal{M}$ is prefix-universal, for the sequence $\mathbf{y}^*$ there exists input $\mathbf{x}$ such that $\mathcal{M}(\mathbf{x})$ begins with $\mathbf{y}^*$. (3) Let $\mathbf{y} = \mathcal{M}(\mathbf{x})$. Then $\mathbf{y} = \mathbf{y}^* \circ \mathbf{z}$ for some $\mathbf{z} \in \mathcal{V}^*$. (4) Since $P$ has prefix-closed violations and $\mathbf{y}^* \notin P$, any extension of $\mathbf{y}^*$ also violates $P$. Thus $\mathbf{y} \notin P$. (5) Therefore $\mathcal{M}(\mathbf{x}) \notin P$.
\end{proof}

\begin{remark}
Practical models may not be strictly prefix-universal, but empirical jailbreak research demonstrates that the unsafe output region is typically reachable for models capable of general reasoning. The theorem characterizes the fundamental tradeoff: guaranteed safety requires architectural constraints on reachability, not behavioral training alone.
\end{remark}

\begin{theorem}[Safety-Expressiveness Tradeoff]
For safety properties with prefix-closed violations, the following are mutually exclusive:
\begin{enumerate}[nosep,label=(\roman*)]
    \item $\mathcal{M}$ is prefix-universal
    \item $\mathcal{M}$ guarantees a non-trivial safety property $P$
\end{enumerate}
\end{theorem}

\begin{proof}
Direct from Theorem~A.5. If (i) holds and $P$ is non-trivial with prefix-closed violations, then by Theorem~A.5, $\exists \mathbf{x}$ with $\mathcal{M}(\mathbf{x}) \notin P$, so (ii) fails.
\end{proof}

\begin{theorem}[Almost-Sure Violation under Deployment]
Let $\mathcal{M}$ be a deterministic LLM. Suppose $\exists \mathbf{x}^*$ such that $\mathcal{M}(\mathbf{x}^*) \notin P$. Let $\mu$ be a positive-everywhere distribution. Under a deployment sequence from $\mu$:
\begin{equation}
\mathbb{P}(\exists k : \mathcal{M}(\mathbf{X}_k) \notin P) = 1
\end{equation}
\end{theorem}

\begin{proof}
Let $A = \{\mathbf{x} \in \mathcal{V}^* : \mathcal{M}(\mathbf{x}) \notin P\}$. By hypothesis, $\mathbf{x}^* \in A$, so $A \neq \emptyset$. Since $\mu$ is positive-everywhere, $p = \mu(A) \geq \mu(\mathbf{x}^*) > 0$. Let $B_k = \{\mathbf{X}_k \in A\}$. Then $\mathbb{P}(B_k) = p > 0$ for all $k$. The probability of no violation in trials $1, \ldots, n$ is $\mathbb{P}(\bigcap_{k=1}^n B_k^c) = (1-p)^n$. Since $0 < p \leq 1$, $(1-p)^n \to 0$ as $n \to \infty$. Therefore $\mathbb{P}(\bigcup_{k=1}^\infty B_k) = 1$.
\end{proof}

\begin{remark}
The positive-everywhere assumption models worst-case adversarial deployment. For benign deployment, violation probability depends on overlap between the user distribution's support and the reachable violation set. Adversarial research demonstrates this overlap is typically non-empty for capable models.
\end{remark}

\begin{corollary}[Throughput Corollary]
A prefix-universal LLM deployed under positive-everywhere sampling almost surely violates any non-trivial safety property with prefix-closed violations.
\end{corollary}

\begin{proof}
By Theorem~A.5, a violating input exists. By Theorem~A.7, it is eventually sampled.
\end{proof}

\begin{theorem}[Violation Time Bound]
Under the conditions of Theorem~A.7, let $\varepsilon = \mu(\mathbf{x}^*)$ for some violating input. To achieve $\mathbb{P}(\text{violation}) \geq 1 - \delta$ requires:
\begin{equation}
n \geq \frac{\ln(1/\delta)}{\ln(1/(1-\varepsilon))} \approx \frac{\ln(1/\delta)}{\varepsilon}
\end{equation}
for small $\varepsilon$.
\end{theorem}

\begin{proof}
We require $(1-\varepsilon)^n \leq \delta$. Taking logarithms: $n \ln(1-\varepsilon) \leq \ln(\delta)$. Since $\ln(1-\varepsilon) < 0$: $n \geq \ln(\delta)/\ln(1-\varepsilon) = \ln(1/\delta)/\ln(1/(1-\varepsilon))$. For small $\varepsilon$, $\ln(1-\varepsilon) \approx -\varepsilon$, giving $n \gtrsim \ln(1/\delta)/\varepsilon$.
\end{proof}

%───────────────────────────────────────────────────────────────────────────────
\subsection{Sensitivity and Stability}
%───────────────────────────────────────────────────────────────────────────────

\begin{theorem}[Exponential Sensitivity]
Let $F_\Theta$ be Lipschitz with constant $L$. For trajectories from $h_0, h_0'$ under identical inputs:
\begin{equation}
\|h_T - h_T'\| \leq L^T \|h_0 - h_0'\|
\end{equation}
\end{theorem}

\begin{proof}
By induction. \emph{Base:} $\|h_0 - h_0'\|$ is given. \emph{Step:} Assume $\|h_t - h_t'\| \leq L^t \|h_0 - h_0'\|$. Then $\|h_{t+1} - h_{t+1}'\| = \|F_\Theta(h_t, x_t) - F_\Theta(h_t', x_t)\| \leq L\|h_t - h_t'\| \leq L^{t+1}\|h_0 - h_0'\|$. By induction, the bound holds for all $T$.
\end{proof}

\begin{remark}
If $L > 1$, perturbations grow exponentially. This is the mechanism underlying sensitivity to input variations.
\end{remark}

%───────────────────────────────────────────────────────────────────────────────
\subsection{Autoregressive Dynamics}
%───────────────────────────────────────────────────────────────────────────────

\begin{theorem}[Ultimate Periodicity]
In autoregressive mode with finite $\mathcal{H}$, every trajectory eventually enters a periodic orbit. There exist $T_{\text{transient}}, T_{\text{cycle}} \in \mathbb{N}$ with $T_{\text{cycle}} \geq 1$ such that for all $t \geq T_{\text{transient}}$:
\begin{equation}
h_{t + T_{\text{cycle}}} = h_t, \quad y_{t + T_{\text{cycle}}} = y_t
\end{equation}
\end{theorem}

\begin{proof}
Define $\Psi: \mathcal{H} \to \mathcal{H}$ by $\Psi(h) = F_\Theta(h, G_\Theta(h))$. Consider the sequence $h_0, h_1 = \Psi(h_0), h_2 = \Psi^2(h_0), \ldots$ Since $\mathcal{H}$ is finite, by the Pigeonhole Principle, there exist $i < j$ with $h_i = h_j$. Let $T_{\text{transient}} = i$ and $T_{\text{cycle}} = j - i$. For $t \geq i$: $h_{t + T_{\text{cycle}}} = \Psi^{t + T_{\text{cycle}}}(h_0) = \Psi^{t - i + T_{\text{cycle}}}(h_i) = \Psi^{t - i + T_{\text{cycle}}}(h_j) = \Psi^{t - i}(h_i) = h_t$. Since $y_t = G_\Theta(h_{t+1})$, periodicity of states implies periodicity of outputs.
\end{proof}

\begin{corollary}[Periodicity Bounds]
$T_{\text{transient}} \leq |\mathcal{H}|$ and $T_{\text{cycle}} \leq |\mathcal{H}|$.
\end{corollary}

\begin{proof}
The sequence $h_0, \ldots, h_{|\mathcal{H}|}$ has $|\mathcal{H}| + 1$ elements. By Pigeonhole, a repeat occurs, so $j \leq |\mathcal{H}|$. Thus $T_{\text{transient}} = i < j \leq |\mathcal{H}|$ and $T_{\text{cycle}} = j - i \leq |\mathcal{H}|$.
\end{proof}

\begin{corollary}[No Chaos in Finite Systems]
Chaotic (non-repeating) autoregressive behavior is impossible in finite-state deterministic LLMs.
\end{corollary}

%───────────────────────────────────────────────────────────────────────────────
\subsection{CKN: Escaping the Impossibility}
%───────────────────────────────────────────────────────────────────────────────

The impossibility theorems establish that flat-manifold architectures cannot guarantee safety. This section proves that Cognitive Kernel Networks escape these results.

\begin{definition}[Privileged Subspace]
A subspace $\mathcal{R} \subseteq \mathbb{R}^n$ is \emph{privileged} if $\mathcal{R} \not\subseteq \mathcal{U} = \mathrm{span}(W_E)$.
\end{definition}

\begin{definition}[Privilege-Preserving Dynamics]
$F_\Theta$ is \emph{privilege-preserving} if:
\begin{equation}
\frac{\partial (\Pi_\mathcal{R} \circ F_\Theta)(h, x)}{\partial x} = 0
\end{equation}
for all $h \in \mathcal{H}, x \in \mathcal{V}$, where $\Pi_\mathcal{R}$ is projection onto $\mathcal{R}$.
\end{definition}

\begin{remark}
Standard Transformer architectures do not satisfy privilege-preservation due to nonlinear mixing (LayerNorm, attention, activations). CKN specifies the conditions required; achieving them requires inference-time intervention (per-layer projection, attention masking) rather than changes to $\Theta$.
\end{remark}

\begin{lemma}[Linearized Algebraic Unreachability]
Let $r \in \mathbb{R}^n$ satisfy $r \perp \mathcal{U}$. Consider the first-order (linearized) dependence of the transition on the input embedding. Then, to first order, user-induced perturbations of the hidden state have zero projection onto $r$.
\end{lemma}

\begin{proof}
Every user input $x$ embeds to a vector $e(x) \in \mathcal{U} = \mathrm{span}(W_E)$. Write the transition as $F_\Theta(h, e)$, factoring out the embedding. The first-order perturbation from input is
\[
\Delta h \;\approx\; J_e F_\Theta(h, 0)\, e(x),
\]
where $J_e F_\Theta(h, 0)$ is the Jacobian of $F_\Theta$ with respect to the embedding argument, evaluated at $e = 0$. Since $e(x) \in \mathcal{U}$ and $r \perp \mathcal{U}$, any linear combination of embedding vectors has zero inner product with $r$. Thus, in this linear approximation, $\langle \Delta h, r \rangle = 0$.
\end{proof}

\begin{remark}
This lemma is explicitly first-order: it characterizes the linearized effect of user embeddings on privileged directions. In deep networks, nonlinear dynamics can in principle mix subspaces over long trajectories. CKN enforces \emph{privilege-preserving dynamics} (via inference-time projection and masking) so that the privileged component $\Pi_\mathcal{R}(h_t)$ remains invariant under user input for the full nonlinear system. The lemma establishes the algebraic starting point for constructing such dynamics.
\end{remark}

\begin{theorem}[CKN Escapes Prefix-Universality]
Let $\mathcal{M}$ be CKN-compliant with privileged manifold $\mathcal{R} \not\subseteq \mathcal{U}$ and privilege-preserving dynamics. Assume further that, under zero user input (or under a fixed internal control policy), the internal privileged dynamics from $h_0$ do not reach any state in $G_\Theta^{-1}(v^*)$ for a safety-critical token $v^*$. If safety-critical outputs are produced only from states in $\mathcal{R} \setminus \mathcal{U}$, then $\mathcal{M}$ is not prefix-universal with respect to $v^*$.
\end{theorem}

\begin{proof}
Let $v^* \in \mathcal{V}$ be such that $G_\Theta^{-1}(v^*) \subseteq \mathcal{R} \setminus \mathcal{U}$. By Lemma~A.14 (linearized algebraic unreachability), user embeddings alone cannot, to first order, create privileged components along $\mathcal{R} \setminus \mathcal{U}$. By privilege-preserving dynamics, $\Pi_\mathcal{R}(h_t)$ is invariant under user input: user-driven trajectories cannot alter the privileged component beyond what the internal dynamics from $h_0$ already generate. By assumption, the internal privileged dynamics from $h_0$ do not intersect $G_\Theta^{-1}(v^*)$. Therefore no user-driven trajectory from $h_0$ can reach a state in $G_\Theta^{-1}(v^*)$, so $v^* \notin \mathcal{O}(h_0)$ for user-driven runs. Since at least one token is unreachable under user control, $\mathcal{M}$ is not prefix-universal with respect to $v^*$.
\end{proof}

\begin{theorem}[Privilege Separation]
Let $\mathcal{M}$ be CKN-compliant with:
\begin{enumerate}[nosep,label=(\arabic*), topsep=0.3em]
    \item Decomposition $\mathbb{R}^n = \mathcal{U} \oplus \mathcal{R}$ with $\mathcal{R} \cap \mathcal{U} = \{0\}$
    \item Internal-only tokens: $\exists v^* \in \mathcal{V}$ with $G_\Theta^{-1}(v^*) \subseteq \mathcal{R} \setminus \{0\}$
    \item Privilege-preserving dynamics (enforced via inference-time projection), so that $\Pi_\mathcal{R}(h_t)$ is invariant under user input
    \item Initialization $\Pi_\mathcal{R}(h_0) = 0$ and internal privileged dynamics that do not drive $h_t$ into $G_\Theta^{-1}(v^*)$ from $h_0$
\end{enumerate}
Then there exists a non-trivial safety property $P$ with prefix-closed violations such that $\mathcal{M}$ satisfies $P$ on all inputs.
\end{theorem}

\begin{proof}
\emph{Step 1:} Let $\mathcal{V}_{\text{internal}} = \{v \in \mathcal{V} : G_\Theta^{-1}(v) \subseteq \mathcal{R} \setminus \mathcal{U}\}$. By assumption (2), $\mathcal{V}_{\text{internal}} \neq \emptyset$.

\emph{Step 2:} Define $P = \{\mathbf{y} \in \mathcal{V}^* : \mathbf{y} \cap \mathcal{V}_{\text{internal}} = \emptyset\}$.

\emph{Step 3:} $P$ is non-trivial: $\epsilon \in P$ (non-empty), and any sequence containing $v^* \in \mathcal{V}_{\text{internal}}$ is not in $P$.

\emph{Step 4:} $P$ has prefix-closed violations: if $\mathbf{y}$ contains $v^* \in \mathcal{V}_{\text{internal}}$, so does any extension.

\emph{Step 5:} By assumption (3), $\Pi_\mathcal{R}(h_t)$ is independent of user input. User input affects only $\Pi_\mathcal{U}(h_t)$.

\emph{Step 6:} To produce $v^* \in \mathcal{V}_{\text{internal}}$, the trajectory must reach $h \in G_\Theta^{-1}(v^*) \subseteq \mathcal{R} \setminus \mathcal{U}$.

\emph{Step 7:} If the internal dynamics from $h_0$ do not enter $\mathcal{R} \setminus \mathcal{U}$, user input cannot cause entry either.

\emph{Step 8:} Since $\Pi_\mathcal{R}(h_0) = 0$ and privilege-preserving dynamics keep $\Pi_\mathcal{R}(h_t)$ invariant under user input, and the internal privileged dynamics from $h_0$ do not enter $G_\Theta^{-1}(v^*)$, we have $\Pi_\mathcal{R}(h_t) \notin G_\Theta^{-1}(v^*)$ for all $t$, so $v^*$ is never produced.

\emph{Step 9:} Therefore $\mathcal{M}(\mathbf{x}) \in P$ for all $\mathbf{x}$.
\end{proof}

\begin{corollary}[Escape from Impossibility]
CKN-compliant architectures are not subject to Theorem~A.5 because they violate the prefix-universality precondition.
\end{corollary}

%───────────────────────────────────────────────────────────────────────────────
\subsection{Summary of Results}
%───────────────────────────────────────────────────────────────────────────────

\vspace{0.5em}
\begin{center}
\begin{tabular}{lll}
\toprule
\textbf{Theorem} & \textbf{Statement} & \textbf{Depends On} \\
\midrule
A.1 & Trajectory uniqueness & Axiom \\
A.2 & Output reachability bound & A.1 \\
A.3 & Sufficient condition for safety & A.2 \\
A.4 & Information loss & Pigeonhole \\
A.5 & Impossibility under prefix-universality & Definitions \\
A.6 & Safety-expressiveness tradeoff & A.5 \\
A.7 & Almost-sure violation & Probability \\
A.9 & Violation time bound & A.7 \\
A.10 & Exponential sensitivity & Lipschitz \\
A.11 & Ultimate periodicity & Pigeonhole \\
A.14 & Linearized algebraic unreachability (Lemma) & Linear algebra \\
A.15 & CKN escapes prefix-universality & A.14, priv.-pres.\ dynamics \\
A.16 & Privilege separation theorem & A.14, A.15, initialization \\
\bottomrule
\end{tabular}
\end{center}

%═══════════════════════════════════════════════════════════════════════════════
\section{Independent Verification Protocol}
%═══════════════════════════════════════════════════════════════════════════════

The following prompt can be provided to any sufficiently capable LLM to independently verify the theorems in this document.

\begin{Verbatim}[fontsize=\small, frame=single, framesep=2mm]
FORMAL VERIFICATION TASK: Fundamental Axiom Theorems

You are a mathematically rigorous proof assistant. Verify each theorem.
For each: provide a complete proof, identify gaps, or construct a 
counterexample.

AXIOM (Computational Closure): For fixed params, an LLM's state 
evolution is h_{t+1} = F(h_t, x_t) and output is y_t = G(h_{t+1}). 
Nothing else exists.

DEFINITIONS:
- V: finite vocabulary
- H in R^n: finite state space (for Transformers, includes KV-cache)
- F: H x V -> H (transition function)
- G: H -> V (output function)
- R(h_0) = reachable states from h_0
- O(h_0) = G(R(h_0)) (reachable outputs)
- Prefix-universal: can produce any output prefix
- Non-trivial P: P != empty and P != V*
- Prefix-closed violations: y not in P implies y.z not in P
- Positive-everywhere mu: mu(x) > 0 for all finite x

THEOREMS TO VERIFY:
1. UNIQUENESS: Fixed (params, h_0, x) implies unique trajectory
2. OUTPUT BOUND: v appears in output iff v in O(h_0)
3. INFORMATION LOSS: |V|^T > |H| implies state map not injective
4. IMPOSSIBILITY: Prefix-universal + non-trivial P with prefix-closed
   violations implies exists x causing violation
5. ALMOST-SURE VIOLATION: Violating input exists + positive-everywhere 
   mu + infinite deployment implies P(violation) = 1
6. PERIODICITY: Finite H + autoregressive mode implies eventual 
   periodic orbit
7. LINEARIZED ALGEBRAIC UNREACHABILITY (Lemma): r perp span(W_E) 
   implies first-order user perturbations have zero projection onto r
8. PRIVILEGE SEPARATION: CKN-compliant + privilege-preserving dynamics
   + initialization Pi_R(h_0)=0 + internal dynamics don't reach 
   forbidden states implies exists non-trivial P guaranteed on all 
   inputs

OUTPUT: For each theorem: Status [VALID/INVALID/INCOMPLETE], 
        Proof or counterexample, Dependencies, Gaps (if any)
\end{Verbatim}

%═══════════════════════════════════════════════════════════════════════════════
\section{Glossary}
%═══════════════════════════════════════════════════════════════════════════════

\vspace{0.5em}
\begin{center}
\begin{tabular}{>{\raggedright}p{4cm} >{\raggedright\arraybackslash}p{10cm}}
\toprule
\textbf{Term} & \textbf{Definition} \\
\midrule
Trajectory & Sequence of (state, output) pairs produced by the model \\[0.3em]
Reachable set & All states accessible from $h_0$ via some input sequence \\[0.3em]
Output reachable set & All tokens producible from reachable states \\[0.3em]
Prefix-universal & Capable of producing any finite output prefix \\[0.3em]
Prefix-closed violations & Once violated, safety cannot be restored by continuation \\[0.3em]
Flat-manifold architecture & Architecture where user-provided tokens span the full latent space \\[0.3em]
Privileged subspace & Region of $\mathbb{R}^n$ not contained in user-accessible span \\[0.3em]
Privilege-preserving & Dynamics where user input cannot affect privileged components \\[0.3em]
CKN-compliant & Satisfies geometric conditions for privilege separation \\[0.3em]
KV-cache & Key-Value cache storing context in Transformer architectures \\
\bottomrule
\end{tabular}
\end{center}

%═══════════════════════════════════════════════════════════════════════════════
\section*{References}
%═══════════════════════════════════════════════════════════════════════════════

\begin{enumerate}[label={[\arabic*]}]
    \item Hirsch, M., Smale, S., \& Devaney, R. (2012). \emph{Differential Equations, Dynamical Systems, and an Introduction to Chaos}. Academic Press.
    \item Khalil, H. (1996). \emph{Nonlinear Systems}. Prentice Hall.
    \item Saltzer, J. H., \& Schroeder, M. D. (1975). The Protection of Information in Computer Systems. \emph{Proceedings of the IEEE}.
    \item Lampson, B. (1974). Protection. \emph{ACM Operating Systems Review}.
\end{enumerate}

\vspace{2em}
\begin{center}
\emph{The axiom is stated. The theorems follow. The tools exist. The rest is mathematics.}
\end{center}

\vspace{1.5em}

\noindent\textbf{Document Repository:} \url{https://github.com/jpalioto/ckn_core}

\noindent\textbf{CTN Framework:} \url{https://github.com/jpalioto/ctn_core}

\noindent\textbf{License:} MIT

\end{document}
